\chapter{讨论}\label{Chap:chap5}

本章讨论Swift算法的设计考量、理论分析、适用边界以及与相关工作的比较。

\section{设计考量}

\subsection{11维观测空间的设计原理}

Swift的11维有效观测子空间面向当前控制规则所需的信息设计，包括ECN状态、拥塞控制状态机状态、拥塞事件类型、RTT、窗口和在途字节。另有4项元数据通道，合计构成15元素跨进程状态容器。各字段由ns-3 TCP拥塞控制回调或仿真环境提供，无需额外探测流量。该状态集合用于当前三类拥塞语义判定，不代表覆盖协议栈的全部可观测状态。由于维度与窗口容量固定，单次处理相对于仿真时长为常数级；本文不报告未经测量的具体决策耗时。

\subsection{差异化窗口保留因子的选择}

Swift的差异化窗口保留因子（丢包0.70、ECN 0.75、超时0.50）的选择基于以下分析：

\textbf{丢包保留因子0.70}：相对于保留因子0.50，0.70减少了单次丢包后的窗口回退幅度。在假设后续采用相同线性增速的简化模型中，恢复缺口由$0.5W$降至$0.3W$，对应恢复步数减少40\%；该比例不是当前实验直接测得的恢复时间改善。

\textbf{ECN保留因子0.75}：ECN标记是队列达到标记条件后的显式信号。保留因子0.75使窗口较温和地缩减，并配合降窗后冻结抑制立即反弹。该取值属于设计参数；第~\ref{Chap:chap4}章的FlowMonitor聚合指标无法分离ECN与丢包响应的独立贡献，仍需消融实验验证。

\textbf{超时保留因子0.50}：CA\_LOSS表示协议栈进入重传超时相关的丢失状态。当前实现采用0.50保留因子并重新进入慢启动，这是三类响应中幅度最大的窗口缩减。

此外，连续缩减保护机制（最多3次连续缩减后保持窗口）的设计基于以下分析：在连续拥塞事件中，窗口经过3次保留因子为0.75的缩减后，约保留原窗口的$0.75^3 \approx 42.2\%$，已经完成显著退让。继续缩减将导致窗口过小，恢复时间过长，反而不利于整体吞吐量。

\subsection{自适应参数调优的设计}

Swift的自适应参数调优机制遵循多因素综合（最小RTT感知的RTT膨胀阈值、性能反馈值EMA趋势、连续增长趋势三个维度提高鲁棒性）、渐进调整（避免剧烈变化）和边界约束（$\alpha \in [0.85, 1.30]$，基准值1.10）三个原则。

\section{理论分析}

\subsection{收敛性分析}

本节从窗口动态、拥塞响应和自适应参数边界三个方面定性讨论稳定性。在无拥塞情况下，窗口接近BDP目标后增长幅度减小；发生拥塞时，差异化保留因子与连续缩减保护限制回退过程；$\alpha$的调整受边界和小步长约束。上述性质说明状态与控制量有界，但当前工作未给出严格收敛证明，也未证明$\alpha$必然收敛到全局最优或唯一取值。

\subsection{公平性分析}

Swift的公平性取决于每流独立状态管理（避免流间干扰）、差异化响应一致性（所有Swift流对相同拥塞信号采用相同保留因子）和与传统算法共存的兼容性。

\subsection{复杂度分析}

\textbf{时间复杂度}：单次决策的时间复杂度为$O(1)$，包括观测空间处理、拥塞检测、参数调优和窗口计算。

\textbf{空间复杂度}：每流状态由常量规模字段构成，包括两个长度40的滑动窗口（投递速率过滤窗口与RTT采样窗口）、一个容量上限4096的投递速率采样队列（其实际占用随路径RTT与确认密度自适应收缩）、BDP估计值、$\alpha$以及若干计数器，因此单流为$O(1)$，总空间复杂度为$O(n)$（$n$为并发流数量）。

\textbf{通信开销}：\texttt{GetSsThresh()}与\texttt{IncreaseWindow()}回调各触发一次同步观测—动作交互。15个64位观测元素的原始载荷为120字节，实际消息还包含Protocol Buffers与ZeroMQ封装开销；当前日志未提供端到端通信时延或总字节数测量。

\section{适用边界与改进方向}
\label{sec:limitation-future}

\subsection{在途水位与延迟的剩余代价}

第~\ref{Chap:chap4}章的结果显示，纯TCP设置下Swift相对NewReno/CUBIC均值的时延最高不超过+1.7\%，相对延迟最低基线最高约+10.4\%（S14 \texttt{wifi\_legacy}，71.22\,ms对64.51\,ms；S5 \texttt{congested\_heavy}，0.956\,ms对NewReno 0.874\,ms；S19，381.3\,ms对367.2\,ms）。UDP突发下的最大相对差异约12.2\%，出现在\texttt{congested\_medium}。这些场景多为低速率、深缓冲或受突发流量影响的链路；较高在途水位可能解释部分额外排队，但当前聚合数据不足以建立单一因果关系。

这些场景的结构特征可为后续诊断提供线索。仿真器按$Q = \max(BDP \times n / MTU, 100)$个数据包为瓶颈队列定尺，在低带宽、短传播时延路径上，队列容量可能由100包下界决定。以S8为例，按瓶颈100\,Mbps和基线RTT 18\,$\mu$s计算的路径BDP约为225字节，而队列容量为150\,kB。Swift另将拥塞窗口限制在$[4 \times MSS, \max(4 \times BDP, 200 \times MSS)]$内；其中$200\times MSS$是上界的兜底项，并不会直接把窗口推至该值。现有FlowMonitor数据没有记录cwnd和队列长度轨迹，因此无法判定时延差异究竟由窗口下界、目标窗口、队列下界还是瞬态过程主导。

后续可通过记录cwnd、在途字节和队列占用轨迹来区分上述因素，并分别评估：（1）BDP估计建立前后的窗口上界切换；（2）以$q_{est}=RTT_{smoothed}-RTT_{min}$约束窗口增长；（3）多时间尺度BDP估计。上述方案均属待验证设计，不作为当前实现的既有能力。

\subsection{长RTT路径的验证与残余问题}

旧版实现采用逐ACK公式$DR = segmentsAcked \times segSize / lastRtt$估计投递速率，该公式会把单次确认载荷除以完整RTT，不能直接表示窗口尺度上的交付率。v0.1.0已将其替换为时间窗口差分形式$DR = [B(t) - B(t-\Delta)]/\Delta$（$\Delta \ge 2 \times RTT_{min}$，见第~\ref{Chap:chap3}章第3.7.1节）。当前数据中，跨域WAN（S9）、长途WAN、LEO与GEO卫星等长RTT场景均呈现吞吐优势，未再观察到旧版数据中的明显低吞吐特征。由于没有保留同一配置下的新旧估计器消融结果，这一变化只能作为与修正方向一致的证据，不能单独证明因果。长RTT路径的增益量级仍需通过多种子和参数扫描进一步评估。

\subsection{实验证据的范围限制}

本文的定量结论受以下范围限制约束，这些限制同时约束机理解释与结果外推。

第一，\textbf{单一随机运行号}。每个(场景, 协议, 设置)三元组仅对应一次仿真运行，因此本文不报告跨运行号的标准差与置信区间，并据此仅把量级明确、机理可解释的差异写入结论。ns-3.40在给定\texttt{RngRun}下确定执行；结果对随机实现的稳健性仍需通过多个\texttt{RngRun}重复确认。

第二，\textbf{队列与流量配置边界}。本章数据全部在RED/ECN对称标记、突发平均负载=瓶颈32\%的配置下生成。该配置是当前实现（v0.1.0）的参考配置，但并不能分离ECN与丢包分支、差异化保留因子与目标窗口机制各自的独立贡献（需要逐机制消融），也不能外推到丢尾/CoDel队列或其他突发强度。完成机制消融与配置扫描是后续工作中优先级最高的一项。

\subsection{参数自动调优}

当前Swift的保留因子（0.70/0.75/0.50）、$\alpha$边界（$[0.85,1.30]$）、连续递减阈值（$D_{max}=3$）和降窗后冻结窗口增长的ACK事件数（4）为当前实现中的手动调优预设值。后续研究可通过元学习\upcite{meta_learning}或贝叶斯优化探索参数的在线自动搜索，但相关结论需要在完整实验矩阵重跑后再报告。

\subsection{学习型策略增强}

Swift当前采用基于规则的启发式策略，其可解释性强但表示能力受人工规则结构限制。若后续引入学习型组件（如离线强化学习训练神经网络策略、再经模型蒸馏提取为紧凑规则），需要引入安全约束并在完整实验矩阵与异常过滤流程上重新评估；在获得此类结果之前，本文不将其作为当前系统的组成部分。

\subsection{真实部署验证}

需将算法实现为Linux内核TCP模块或QUIC用户态协议栈，在覆盖近端高速、蜂窝/无线接入、广域网与卫星链路的真实或半实物测试床上验证。仿真与真实环境的差异（如NIC offload、中断合并、NUMA效应、硬件时间戳精度、终端操作系统调度和无线链路波动）可能影响实际性能。QUIC协议的用户态实现可简化部署流程，且QUIC的多流特性可实现更精细的流间调度。

\section{与相关工作的比较}

\subsection{与Aurora的比较}

Aurora\upcite{aurora}是基于深度强化学习的拥塞控制算法，与不采用学习型策略的Swift的主要区别包括：

\begin{table}[htbp]
\centering
\caption{Swift与Aurora的详细对比}
\label{tab:vs_aurora}
\begin{tabular}{|l|c|c|}
\hline
\textbf{特性} & \textbf{Swift} & \textbf{Aurora} \\
\hline
有效观测维度 & 11（另含4项元数据） & 10 \\
ECN状态利用 & 是 & 否 \\
CA状态利用 & 是 & 否 \\
决策机制 & 基于规则 & 深度神经网络 \\
可解释性 & 高 & 低 \\
训练需求 & 无 & 需要大量训练 \\
决策复杂度 & $O(1)$ & 神经网络前向推理 \\
参数调优 & 在线自适应 & 固定（训练后） \\
\hline
\end{tabular}
\end{table}

\subsection{与DCTCP的比较}

DCTCP\upcite{dctcp}是利用ECN进行细粒度调整的拥塞控制算法。

\begin{table}[htbp]
\centering
\caption{Swift与DCTCP的详细对比}
\label{tab:vs_dctcp}
\begin{tabular}{|l|c|c|}
\hline
\textbf{特性} & \textbf{Swift} & \textbf{DCTCP} \\
\hline
ECN利用方式 & 状态直接检测 & 标记比例连续调整 \\
拥塞响应 & 差异化保留因子 & 统一公式 \\
丢包响应保留因子 & 0.70 & 0.50 \\
ECN响应保留因子 & 0.75 & 可变（基于$\alpha$） \\
参数自适应 & 是 & 否 \\
协议栈状态利用 & 是 & 否 \\
多信号融合 & 是 & 否 \\
\hline
\end{tabular}
\end{table}

\subsection{与BBR的比较}

BBR\upcite{bbr}是基于带宽和延迟估计的拥塞控制算法。

\begin{table}[htbp]
\centering
\caption{Swift与BBR的详细对比}
\label{tab:vs_bbr}
\begin{tabular}{|l|c|c|}
\hline
\textbf{特性} & \textbf{Swift} & \textbf{BBR} \\
\hline
主要拥塞信号 & 丢包+ECN+CA状态 & RTT测量 \\
带宽探测方式 & 自适应参数调优 & 周期性增益调整 \\
丢包响应 & 是 & 否（设计上） \\
ECN利用 & 是 & 否 \\
控制模式 & 窗口控制 & 速率控制 \\
低RTT场景带宽利用率 & 接近饱和 & 场景相关 \\
中速场景带宽利用率 & 较高 & 较高 \\
与CUBIC共存公平性 & 场景相关，依赖队列与RTT条件 & 场景相关，依赖增益周期 \\
\hline
\end{tabular}
\end{table}

\subsection{与PCC的比较}

PCC\upcite{pcc}采用在线实验探索最优速率，而Swift采用预设规则结合自适应参数调优。PCC需要在线进行速率探索实验，可能导致短期性能下降；Swift不需要显式探索，性能更稳定。

\section{敏感性分析}

参数敏感性的理论分析应服务于对实现机制的解释。当前实现采用固定边界与事件驱动更新逻辑：ECN保留因子为0.75，普通丢包保留因子为0.70，超时保留因子为0.50；乘性增加因子$\alpha$在$[0.85,1.30]$范围内根据RTT膨胀、性能反馈值EMA和连续增长趋势动态调整。本节不引入未由日志支撑的额外参数扫描结果，而从机制角度说明这些取值如何影响稳定性与吞吐量。

\subsection{窗口保留因子的理论依据}

丢包保留因子0.70的选择基于收敛速度与吞吐量的权衡：在AIMD框架下，保留因子$\beta$决定了窗口从$W_{max}$恢复到$W_{max}$所需的时间$T_{recover} \propto \frac{1-\beta}{\gamma}$个RTT。$\beta=0.70$相比传统$\beta=0.50$使恢复时间缩短40\%，但丢包后的即时退让幅度也相应减小。在有ECN支持的网络中，丢包事件频率显著降低，因此更温和的丢包响应是合理的。

ECN保留因子0.75的选择基于ECN信号的语义：ECN标记表示队列深度超过阈值但尚未溢出，属于早期拥塞信号。相较普通丢包响应，0.75保留因子能够在发生早期拥塞通知时更温和地释放队列压力；相较超时响应，它又避免了不必要的窗口深度回退。该参数与降窗后冻结机制共同作用，用于降低连续ACK事件中立即反弹造成的二次排队风险。

\subsection{连续递减保护阈值的数学分析}

在连续ECN标记事件中，窗口经过$D_{max}$次保留因子为$\beta_{ecn}=0.75$的收缩后，剩余窗口比例为$0.75^{D_{max}}$。当$D_{max}=3$时，算法在连续拥塞信号下已经完成显著退让；继续无约束缩减容易造成窗口过低和恢复迟滞。因此实现中采用连续递减保护：当连续缩减计数超过阈值后暂时保持当前窗口，同时通过降窗后冻结若干ACK事件抑制立即增窗，使队列有时间排空。

\subsection{RTT膨胀比阈值的设计考量}

RTT膨胀调节采用路径感知的动态阈值，而非固定阈值。实现首先根据最小RTT构造松弛量$s = 0.3 + 0.15 \sqrt{\max(0, RTT_{min,ms}-1)}$，再形成$1+s$、$1+2s$和$1+4s$三个阈值。短RTT路径对排队更敏感，阈值相对更紧；长RTT路径允许更大的相对抖动，阈值相对放宽。该设计避免在远距离路径上因正常RTT波动过早收紧窗口，也避免在低RTT路径上放任队列持续积累。

\section{计算开销分析}
\label{sec:overhead}

Swift当前实现采用事件驱动的规则决策与每流独立状态表。由于本仓库日志未包含跨硬件平台的微基准测量，本节不报告未验证的决策延迟、内存占用或CPU利用率数值，而从代码结构给出可复核的复杂度分析。

\subsection{时间复杂度}

每次ACK或拥塞事件触发时，算法只执行常数次状态读取、拥塞类型判定、$\alpha$更新和窗口计算。带宽估计使用固定长度为40的滑动窗口，RTT采样窗口长度同样为40，且最大交付速率由窗口队列维护。因此，单次决策的时间复杂度为$O(1)$，不会随仿真总时长或历史样本数量线性增长。

\subsection{空间复杂度}

每个TCP连接维护独立状态，包括长度为40的交付速率窗口、长度为40的RTT采样窗口、BDP估计值、$\alpha$、连续增长/缩减计数器、丢包/ECN计数器和降窗后冻结计数器等常量规模字段。因此，单流空间复杂度为$O(1)$，总空间复杂度随并发连接数线性增长，即$O(N)$。

\subsection{工程开销来源}

主要开销来自ZeroMQ同步状态—决策交互、Python规则决策和ns-3回调状态采集。当前实现优先保证跨进程决策与离散事件仿真时序一致，因此未在同步信道上启用MTP多线程并行仿真。若未来进行内核态或用户态协议栈移植，应重新测量实际平台上的决策时延和资源占用。

\section{理论分析深入探讨}

\subsection{多信号融合的信息论基础}

多信号融合拥塞检测的设计可以从信息论角度进行分析。设网络的真实拥塞状态为随机变量$C \in \{0, 1\}$，观测到的信号集合为$\mathbf{S} = \{S_1, S_2, ..., S_k\}$。根据贝叶斯定理：
\begin{equation}
P(C=1|\mathbf{S}) = \frac{P(\mathbf{S}|C=1) \cdot P(C=1)}{P(\mathbf{S})}
\end{equation}

Swift的分层优先级机制可以理解为对不同拥塞信号可靠性的排序：慢启动阈值回调对应的丢包/超时信号具有最高优先级；ECN CE或ECE状态表示队列已经触发显式拥塞通知，具有次级优先级；RTT膨胀和恢复状态只参与参数调节，不作为独立降窗触发条件。这样的排序与实现中的\texttt{\_detect\_congestion}函数保持一致。

\subsection{差异化响应的最优控制分析}

差异化窗口保留因子的设计可以从最优控制理论角度分析。优化目标是最大化长期吞吐量同时最小化延迟：
\begin{equation}
J = \sum_{t=0}^{T} \left[ -\alpha \cdot throughput(t) + \beta \cdot delay(t) + \gamma \cdot |u(t)|^2 \right]
\end{equation}

Swift的差异化保留因子（丢包0.70、ECN 0.75、超时0.50）可以理解为对不同拥塞严重程度的分层近似。

\subsection{自适应参数调优的规则化描述}

Swift当前的参数自适应规则可以表示为如下分段函数：
\begin{equation}
\pi(s_t) = \begin{cases}
+\Delta_{\text{up}} & \text{如果~} r < 1+s \text{，表示排队膨胀较低} \\
+\Delta_{\text{small}} & \text{如果~} 1+s \le r < 1+2s \text{，表示温和增长} \\
-\Delta_{\text{down}} & \text{如果~} r > 1+4s \text{，表示重度排队} \\
+\Delta_{\text{ema}} & \text{如果~} R_{fast} > R_{base} + m \\
-\Delta_{\text{ema}} & \text{如果~} R_{fast} < R_{base} - 4m \\
+\Delta_{\text{stable}} & \text{如果~} consecutive\_growth > 8 \\
0 & \text{否则}
\end{cases}
\end{equation}
其中$s = 0.3 + 0.15\sqrt{\max(0, RTT_{min,ms}-1)}$，$r=RTT_{last}/RTT_{min}$，$R_{fast}$与$R_{base}$分别为性能反馈值的快线（$\eta_f=0.15$）与慢线基线（$\eta_b=0.02$）指数移动平均，$m = \max(0.25, 0.1|R_{base}|)$为自适应容差。该形式体现了当前实现中的两项关键设计：路径感知的动态RTT阈值，以及以基线相对偏离而非绝对水平判断反馈趋势——后者的必要性在于环境侧反馈值在绝大多数ACK事件上为正，绝对阈值会使调节退化为单向棘轮。

\subsection{收敛性的Lyapunov分析}

定义Lyapunov函数：
\begin{equation}
V(t) = \frac{1}{2}(cwnd(t) - cwnd^*)^2 + \frac{1}{2}(queue(t) - queue^*)^2
\end{equation}

当$cwnd > cwnd^*$且出现ECN拥塞通知时，当前实现按$cwnd(t+1) = 0.75 \cdot cwnd(t)$进行窗口退让；若出现普通丢包或超时，则分别按0.70或0.50执行更强响应。该分层缩减机制有助于使窗口向目标区间回落。

\subsection{公平性的定性分析}

多流公平性同时受拥塞控制规则、流启动时刻、RTT、队列规则和丢包/标记分布影响。Swift为各连接维护独立状态，并对相同类型的拥塞信号使用相同响应规则，这为同一路径上的一致处理提供了基础，但不能推出与其他算法竞争时的理论均衡。

第~\ref{Chap:chap4}章中，Swift在纯TCP设置下的Jain指数均值为0.9989、最小值为0.9967，四种协议的场景均值均接近1。逐场景看，Swift在11个场景严格最优，在32个场景中与最优基线的差值不超过0.002；UDP突发下对应数量为17个和35个。该结果支持"总体与基线相当"的描述，不支持逐场景最优或由BDP目标必然产生公平性的因果结论。

\section{本章小结}

本章围绕三项核心技术讨论了Swift的设计边界：多信号拥塞状态感知为超时、ECN驱动的窗口缩减和普通丢包提供语义区分；时间窗口交付率、最小RTT感知阈值与基线相对反馈共同构成启发式自适应控制；差异化保留因子、窗口边界、连续缩减保护、降窗后冻结和陈旧动作作废用于限制窗口过度变化。

现有证据表明，Swift在所考察场景中多数呈现约2\%--6\%的吞吐优势，但时延优势不普适：纯TCP下相对BBR平均高约0.8\%，个别场景最高约10.4\%；UDP突发下相对最低时延基线的最大差异约12.2\%。公平性总体接近1且与基线相当，丢包率也不存在逐场景一致优势。当前结果来自单一随机种子和RED/ECN配置，且缺少机制消融、窗口轨迹与真实终端测试，因此多种子重复、配置扫描和真实协议栈验证仍是必要工作。
